\section*{Bab \ref{ch:b2:ore-deposits} --- Endapan Bijih}

\begin{solution}{exo:b2:ore-deposits:1}
(a) $\qty{0.5}{\percent} = \qty{5000}{ppm}$, dan $5000/28 = 179$. (b) $\qty{6}{g/t} = \qty{6000}{ng/g}$,
dan $6000/1.5 = 4000$. (c) Kromit memuat $2 \times 52.00/(55.85 + 104.00 + 64.00) = 0.4646$
kromium menurut massa; kromitit, $0.9 \times 0.4646 = 0.418$, atau \qty{418000}{ppm}, dan
$418000/92 = 4540$.
\end{solution}

\begin{solution}{exo:b2:ore-deposits:2}
(a) $914 \times 0.0053 = \qty{4.84}{Mt}$ tembaga dan $914 \times 10^6 \times 0.40 = \num{3.66e8}$ g,
yaitu \qty{366}{t} emas; $1640 \times 0.0044 = \qty{7.22}{Mt}$ dan \qty{738}{t}. (b) $1640 - 914 =
\qty{726}{Mt}$ yang memuat $7.216 - 4.844 = \qty{2.37}{Mt}$ tembaga, dengan kadar rata-rata
$2.372/726 = \qty{0.327}{\percent}$: blok-blok tambahan itu terletak di bawah \omterm{def:b2:ore-deposits:grade}{kadar batas} \qty{0.3}{\percent}
dan tepat di atasnya. (c) Keduanya menjawab pertanyaan itu pada dua \omterm{def:b2:ore-deposits:grade}{kadar batas}; tidak satu pun merupakan ukuran
endapan tanpa \omterm{def:b2:ore-deposits:grade}{kadar batasnya}.
\end{solution}

\begin{solution}{exo:b2:ore-deposits:3}
(a) $(19.3 - 1.00)/(2.65 - 1.00) = 18.3/1.65 = 11.1$. (b) Kecepatan yang sama memerlukan $(\rho_g - 1)d_g^2 =
(\rho_q - 1)d_q^2$, sehingga $d_g/d_q = \sqrt{1.65/18.3} = 0.300$: butir emas mengendap seperti butir kuarsa
yang kira-kira tiga kali lebih besar. (c) $d_q = 0.1\sqrt{3.7/1.65} = \qty{0.150}{mm}$.
\end{solution}

\begin{solution}{exo:b2:ore-deposits:4}
Per satuan luas dan densitas, selimut itu mempertahankan tembaganya sendiri dan menerima tembaga dari kolom
yang terlindi: $g = g_0(h_e + \eta h_l)/h_e$. (a) $0.4 \times (50 + 150)/50 = \qty{1.6}{\percent}$.
(b) $0.4 \times (50 + 0.8 \times 150)/50 = \qty{1.36}{\percent}$. (c) $h_l = 50 \times (1.7/0.4 - 1)
= \qty{163}{m}$. Kadar-kadar ini berada dalam rentang 0.4--\qty{1.7}{\percent} selimut porfiri
yang diperkaya.
\end{solution}

\begin{solution}{exo:b2:ore-deposits:5}
(a) Platina, $C_S = 10 \times 10^4 (R + 1)/(R + 10^4)$ ppb: \qty{1.00}{ppm} untuk $R = 100$,
\qty{9.10}{ppm} untuk \num{1000}, \qty{50.0}{ppm} untuk \num{10000}, \qty{90.9}{ppm} untuk \num{100000}.
Nikel, $C_S = 200 \times 500 (R + 1)/(R + 500)$ ppm: \qty{1.68}{\percent}, \qty{6.67}{\percent},
\qty{9.52}{\percent}, \qty{9.95}{\percent}. (b) $D\,C_L$: \qty{100}{ppm} platina,
\qty{10}{\percent} nikel. (c) Nikel berada dalam \qty{5}{\percent} dari batas atasnya mulai
$R = \num{10000}$, sehingga kandungannya hampir tidak berubah sepanjang faktor sepuluh dalam $R$; platina di sana masih
hampir berlipat dua, sehingga kandungannya mengukur $R$.
\end{solution}

\begin{solution}{exo:b2:ore-deposits:6}
(a) $\mu_M = 2.373 - 0.364 - 2 = 0.009$; $\sigma_M = 0.702$; median $10^{0.009} = \qty{1.02}{Mt}$.
(b) $z = (1 - 0.009)/0.702 = 1.412$, $1 - \Phi(1.412) = 0.0790$: sekitar 33 endapan dari 422.
(c) $\sigma = \sqrt{0.654^2 + 0.199^2} = 0.684$, $z = 1.449$, $1 - \Phi = 0.0737$, sekitar 31: korelasi
yang lemah itu menambahkan beberapa raksasa.
\end{solution}

\begin{solution}{exo:b2:ore-deposits:7}
Rata-rata semua blok $0.40\,e^{0.18} = \qty{0.479}{\percent}$. Untuk $c = 0.3$: $z = \ln(0.40/0.3)/0.60 =
0.4795$, $\Phi(z) = 0.6842$, $\Phi(z + s) = 0.8598$: $T = \qty{684}{Mt}$, $\bar g = 0.479 \times
0.8598/0.6842 = \qty{0.602}{\percent}$, $M = 1000 \times 0.00479 \times 0.8598 = \qty{4.12}{Mt}$.
Untuk $c = 0.6$: $z = -0.6758$, $\Phi = 0.2496$, $\Phi(z + s) = 0.4698$: \qty{250}{Mt} dengan
\qty{0.901}{\percent}, \qty{2.25}{Mt}. Pada gambar, pada \omterm{def:b2:ore-deposits:grade}{kadar batas} \qty{0.3}{\percent} kurva biru
terbaca sekitar \qty{680}{Mt} dan kurva merah sekitar \qty{0.6}{\percent}.
\end{solution}

\begin{solution}{exo:b2:ore-deposits:8}
(a) Fluida: $\num{5e6}/(\num{1e-3} \times 0.5) = \qty{1.00e10}{t}$; \omterm{def:b2:magma-emplacement:magma}{magma}: $\num{1e10}/0.035 =
\qty{2.86e11}{t}$. (b) $\num{2.86e14}/2600 = \qty{1.10e11}{m^3}$, \qty{110}{km^3}, sebuah kubus \omterm{def:b2:magma-emplacement:magma}{magma}
bersisi sekitar \qty{4.8}{km}. (c) \omterm{def:b2:magma-emplacement:magma}{Magma} itu memuat $\num{2.86e11}
\times \num{50e-6} = \qty{1.43e7}{t}$ tembaga, dan bijih memuat $5/14.3 = 0.350$ darinya.
\end{solution}

\begin{solution}{exo:b2:ore-deposits:9}
(a) Rata-rata logaritma 1.411, simpangan baku 0.0664; median $10^{1.411} = \qty{25.8}{ppm}$;
ambang $10^{1.411 + 0.133} = 10^{1.544} = \qty{35.0}{ppm}$. (b) $(4 \times 400 + 36 \times 25.8)/40 =
\qty{63.2}{ppm}$, di atas \qty{35.0}{ppm}: masih anomal setelah diencerkan oleh sembilan bagian sedimen
bersih. (c) $4 \times 400 + (A - 4) \times 25.8 = 35.0\,A$ memberikan $A = 4 \times (400 -
25.8)/(35.0 - 25.8) = \qty{163}{km^2}$.
\end{solution}

\begin{solution}{exo:b2:ore-deposits:10}
(a) $e^{su}\,e^{-u^2/2} = e^{-(u - s)^2/2 + s^2/2}$; bagi dengan $\sqrt{2\pi}$. (b) $E[e^{\mu + su}] =
e^{\mu}\int e^{su}\varphi(u)\,\mathrm{d}u = e^{\mu + s^2/2}\int\varphi(u - s)\,\mathrm{d}u = e^{\mu +
s^2/2}$. (c) $\Phi$ naik dan $s > 0$, sehingga $\Phi(z + s) > \Phi(z)$: blok-blok yang dipertahankan lebih kaya daripada
blok-blok yang dibuang, sehingga setiap \omterm{def:b2:ore-deposits:grade}{kadar batas} mempertahankan bagian logam yang lebih besar daripada bagian batuan; selisihnya
bertambah seiring $s$, yaitu variabilitas endapan itu.
\end{solution}

\begin{solution}{exo:b2:ore-deposits:11}
(a) $D(R + 1)/(R + D) = fD$ memberikan $R + 1 = fR + fD$, $R = (fD - 1)/(1 - f)$. (b) $f = 0.5$: $R = D - 2$,
298 dan \num{9998}; $f = 0.9$: $R = 9D - 10$, \num{2690} dan \num{89990}. (c) Nikel mencapai
\qty{50}{\percent} batas atasnya di dalam rentang bijih nikel--tembaga (100 sampai \num{2000}), dan 90 persen tepat di atasnya;
platina memerlukan $R$ sebesar \num{10000} sampai \num{100000}, rentang terumbu PGE.
\end{solution}

\begin{solution}{exo:b2:ore-deposits:12}
(a) Dengan $\ln M = m + s u$, endapan di atas kuantil itu memiliki $u > \Phi^{-1}(q)$; logamnya adalah
$\int_{\Phi^{-1}(q)}^{\infty} e^{m + su}\varphi(u)\,\mathrm{d}u = e^{m + s^2/2}[1 - \Phi(\Phi^{-1}(q) -
s)]$, dan membaginya dengan total $e^{m + s^2/2}$ memberikan $\Phi(s - \Phi^{-1}(q))$. (b) $s = 0.702 \ln 10 =
1.616$; $\Phi(1.616 - 1.282) = \Phi(0.334) = 0.631$. (c) $\Phi(1.616 - 2.326) = \Phi(-0.710) = 0.239$,
dibandingkan dengan 0.266 untuk keempat endapan terbesar: endapan-endapan terbesar sedikit lebih besar daripada yang diharapkan hukum
log-normal.
\end{solution}

\begin{solution}{pb:b2:ore-deposits:1}
\textbf{1.} \qty{330}{m}; $\sum l\,g = 138.4$, rata-rata $138.4/330 = \qty{0.419}{\percent}$.
\textbf{2.} $3.24/8 = \qty{0.405}{\percent}$; setiap selang mewakili suatu panjang inti, jadi suatu
volume batuan: selang yang panjang harus lebih diperhitungkan daripada selang yang pendek.
\textbf{3.} $230/330 = 0.697$, dengan $118.6/230 = \qty{0.516}{\percent}$.
\textbf{4.} $4190/28 = 150$.
\textbf{5.} $8100/28 = 289$.
\textbf{6.} $0.38\,e^{0.55^2/2} = 0.38\,e^{0.1513} = \qty{0.442}{\percent}$; lubang itu, dengan
\qty{0.419}{\percent}, berada dalam \qty{6}{\percent} darinya: representatif.
\textbf{7.} $\pi \times 450^2 \times 700 = \qty{4.45e8}{m^3}$; $\times 2.7 = \num{1.20e9}$ t,
\qty{1200}{Mt}.
\textbf{8.} $z = \ln(0.38/0.3)/0.55 = 0.4298$; $\Phi(z) = 0.6663$.
\textbf{9.} $z = 1.167$, 0.4298, $-0.4990$; $\Phi = 0.8784$, 0.6663, 0.3089; $T = \qty{1050}{Mt}$,
\qty{800}{Mt}, \qty{371}{Mt}.
\textbf{10.} $\Phi(z + s) = \Phi(0.9798) = 0.8364$; $\bar g = 0.442 \times 0.8364/0.6663 =
\qty{0.555}{\percent}$.
\textbf{11.} $0.442 \times 0.9570/0.8784 = \qty{0.482}{\percent}$; $0.442 \times 0.5203/0.3089 =
\qty{0.744}{\percent}$.
\textbf{12.} \omterm{def:b2:ore-deposits:grade}{Tonase} $\times\, 370.7/799.6 = 0.464$; kadar $\times\, 0.744/0.555 = 1.34$.
\textbf{13.} Pada median, \qty{0.38}{\percent} ($z = 0$): \qty{600}{Mt} dengan $0.442 \times
\Phi(0.55)/0.5 = 0.442 \times 0.7088/0.5 = \qty{0.627}{\percent}$.
\textbf{14.} $1200 \times 0.00442 = \qty{5.30}{Mt}$.
\textbf{15.} $(\log_{10} 1200 - 2.373)/0.654 = (3.0792 - 2.373)/0.654 = 1.080$; $\Phi = 0.8599$:
persentil ke-86, suatu endapan besar.
\textbf{16.} $(\log_{10} 0.442 + 0.364)/0.199 = (-0.3546 + 0.364)/0.199 = 0.047$; $\Phi = 0.5187$: kadar
yang biasa, persentil ke-52.
\textbf{17.} 0.702 dengan korelasi, 0.684 tanpa korelasi.
\textbf{18.} $\mu_M = 0.009$; $(\log_{10} 5.30 - 0.009)/0.702 = (0.7243 - 0.009)/0.702 = 1.019$; $\Phi =
0.8459$: persentil ke-85.
\textbf{19.} $422 \times (1 - 0.8459) = 65$.
\textbf{20.} $M = T_0\,e^{\mu + s^2/2}\,\Phi(z + s) = 1200 \times 0.00442 \times 0.9570 =
\qty{5.08}{Mt}$.
\textbf{21.} $1200 \times 0.00442 \times 0.5203 = \qty{2.76}{Mt}$.
\textbf{22.} $\Phi(z + s) = 0.8364$ dari tembaga, dalam $\Phi(z) = 0.6663$ dari \omterm{def:b2:ore-deposits:grade}{tonase}.
\textbf{23.} Dari 0.3 ke \qty{0.2}{\percent} \omterm{def:b2:ore-deposits:grade}{tonase} bertambah sebesar $1054 - 800 = \qty{254}{Mt}$
(\qty{32}{\percent}) dan tembaga sebesar $5.08 - 4.44 = \qty{0.64}{Mt}$ (\qty{14}{\percent}): blok-blok yang ditambahkan
rata-rata berkadar $0.64/254 = \qty{0.25}{\percent}$, yang termiskin dalam endapan itu.
\textbf{24.} $1200 \times 0.00442 \times 0.8364 = 800 \times 0.00555$: prospek itu memuat
\textbf{\qty[text-series-to-math=true]{4.44}{Mt}} tembaga di atas \omterm{def:b2:ore-deposits:grade}{kadar batas} \qty{0.3}{\percent}.
\end{solution}
